\section{Adoption}\label{sec:adoption}

A construction made in one Bonfire becomes part of another's language by the route a word takes between people.
Someone meets a form they cannot read, learns it from a source that can, and keeps it if it goes on working. The
protocol gives that route five steps; Figure~\ref{fig:adoption} shows the route.

\paragraph{Signal.} Each Bonfire publishes, with its grammar, a signal: a statement of what the grammar is for.
Signals are how a grammar is found. They cost nothing to make, so a Bonfire could claim every purpose; for that
reason no payment rests on a signal, and each Bonfire ranks the grammars it consults by its own record of which
ones answered it.

\paragraph{Lookup.} When a Bonfire's agent meets a form its inventory cannot read, the gap becomes a \emph{lookup}:
a query to the network, routed by the signals into the grammars of other Bonfires. Published grammars are open, so
the lookup runs on a node that serves the asking Bonfire, which fetches the foreign grammar and applies it. The text that
could not be read is shown to nobody else. Lookups that no grammar answers, grouped by what they were about, are
the network's record of what remains to be structured.

\paragraph{Trial.} If a foreign grammar reads the text, the constructions necessary to that reading are on offer
together, since a construction torn from those it depends on may read nothing. They enter the asking Bonfire's
inventory by reference to their originator's registration, at the starting score, as a \emph{trial}. A construction
on trial reads, and its effect is switched off: an effect taken from elsewhere is behaviour taken from elsewhere,
and it should not act before someone in the Bonfire has agreed to it.

\paragraph{Selection.} Over the Bonfire's later discourse a construction on trial succeeds or fails beside the
Bonfire's own, and its standing moves by the rule of Section~\ref{sec:constructions}. Where several Bonfires offer
a reading for the same gap, their constructions compete here, and different Bonfires may settle differently. This
is the naming game, played between communities \citep{steels1995spatial,baronchelli2006sharp}.

\paragraph{Adoption.} The Bonfire's stakeholders confirm a construction that has earned its place, and it stays in the Bonfire's inventory with its effect on. It does not become one of the Bonfire's own
constructions: its originator is unchanged. The confirmation is a signed attestation that references the
originator's registration; that act is the \emph{adoption}. From then on, a Bonfire \emph{holds} a construction in
an epoch when it has adopted it and the construction did necessary work in that Bonfire's own discourse in the
epoch, as shown by receipts that anyone can re-execute. Holding is what the protocol counts
(Section~\ref{sec:reward}).

Adoption is free, and it is pulled by need: no Bonfire can place a construction in another. A Bonfire that starts
from nothing becomes useful by adopting, and it has a reason to adopt by reference instead of copying, since a
reference brings the construction as its originator maintains it and a copy earns its copier nothing
(Section~\ref{sec:lineage}). A group that begins from the whole of another Bonfire's grammar is in this sense not a
rival to it: it is one more community that holds its constructions.

\begin{figure}[!htb]
\centering
% fig:adoption -- a gap in B, a lookup into A's grammar, trial and adoption, and what flows back to A
\begin{tikzpicture}[
  font=\small,
  cxn/.style={draw=cxnnavy, fill=cxngreen, rounded corners=2pt, minimum width=3.0cm, minimum height=0.8cm,
              align=center},
  plain/.style={draw=black!60, rounded corners=2pt, minimum width=3.0cm, minimum height=0.8cm, align=center},
  desc/.style={-{Stealth[length=5pt]}, thick, cxnnavy},
  credit/.style={-{Stealth[length=5pt]}, thick, okgreen, dashed},
  note/.style={font=\scriptsize, text=black!70, align=center}
]
  \node[cxn] (pa) at (0,0) {construction\\[-1pt]{\scriptsize registered by Bonfire $A$}};
  \node[draw=modelred, text=modelred, rounded corners=2pt, minimum height=0.8cm, align=center, font=\scriptsize]
    (gap) at (4.6,0) {gap in\\Bonfire $B$};
  \node[plain] (ad) at (9.2,0) {trial, then adoption\\[-1pt]{\scriptsize by $B$'s stakeholders}};
  \node[cxn] (var) at (9.2,-2.0) {variant\\[-1pt]{\scriptsize registered by Bonfire $B$}};
  \draw[desc, modelred] (gap) -- node[above, font=\scriptsize, text=modelred] {lookup} (pa);
  \draw[desc] (gap) -- node[above, font=\scriptsize, text=black!70] {by reference} (ad);
  \draw[desc, black!60] (ad) -- node[right, font=\scriptsize, text=black!70] {changed} (var);
  \draw[credit] (ad.north) -- ++(0,0.55) -| node[pos=0.25, above, font=\scriptsize, text=okgreen]
    {$B$ counts toward the breadth of $A$'s construction} (pa.north);
  \draw[credit] (var.west) -| node[pos=0.25, above, font=\scriptsize, text=okgreen]
    {parent: share $\lambda$ of the variant's credit} (pa.south);
\end{tikzpicture}

\caption{Adoption. A form from Bonfire $A$ meets a gap in Bonfire $B$. $B$'s lookup finds $A$'s construction, tries
it, and adopts it; $B$ then counts toward the breadth of $A$'s construction for as long as $B$ holds it. If $B$
changes the construction it registers a variant with $A$'s as its parent, and a share $\lambda$ of the variant's
credit passes to $A$.}\label{fig:adoption}
\end{figure}
